\section{Fungsi Lanjutan: Expression, Arrow, Default, dan Rest}

Bab 8 memperkenalkan fungsi dasar; section ini memperdalam fitur fungsi modern ES6+. \textbf{Function expression} menyimpan fungsi anonim di variabel: \code{const greet = function(name) \{ ... \}}. \textbf{Arrow function} (\code{=>}) menyederhanakan sintaks dan memiliki perilaku \code{this} leksikal---ia mewarisi \code{this} dari scope luar, bukan menciptakan \code{this} sendiri \cite{jsInfo}.

\begin{table}[htbp]
\centering
\begin{tabular}{|P{3cm}|P{4.5cm}|P{4.5cm}|}
\hline
\textbf{Fitur} & \textbf{Sintaks} & \textbf{Catatan} \\
\hline
Default parameter & \code{function f(x = 10)} & Nilai default jika argumen undefined \\
\hline
Rest parameter & \code{function f(...args)} & Menampung sisa argumen ke array \\
\hline
Spread (panggil) & \code{f(...arr)} & Membongkar array menjadi argumen \\
\hline
Arrow function & \code{(x) => x * 2} & Tanpa \code{this} sendiri; ringkas \\
\hline
IIFE & \code{(() => \{ ... \})()} & Langsung dieksekusi; isolasi scope \\
\hline
\end{tabular}
\caption{Fitur fungsi lanjutan ES6+}
\end{table}

\begin{konsep}
\textbf{Higher-Order Function.} Fungsi yang menerima fungsi lain sebagai argumen atau mengembalikan fungsi disebut \textit{higher-order function}. Contoh: \code{arr.map(fn)}, \code{arr.filter(fn)}, \code{setTimeout(fn, ms)}. Konsep ini adalah fondasi pemrograman fungsional di JavaScript dan memungkinkan kode yang deklaratif dan modular \cite{eloquentJs}.
\end{konsep}

